\documentclass[11pt]{article}
\usepackage[margin=1in]{geometry}
\usepackage{amsmath,amssymb,amsthm,mathtools}
\usepackage[T1]{fontenc}
\usepackage{lmodern,microtype}
\usepackage[colorlinks=true,linkcolor=blue,urlcolor=blue,citecolor=blue]{hyperref}
\setlength{\emergencystretch}{2em}
\newtheorem{theorem}{Theorem}
\newtheorem{lemma}[theorem]{Lemma}
\newtheorem{proposition}[theorem]{Proposition}
\theoremstyle{definition}
\newtheorem{definition}[theorem]{Definition}
\newtheorem{remark}[theorem]{Remark}
\newcommand{\ket}[1]{\lvert #1\rangle}
\newcommand{\norm}[1]{\lVert #1\rVert}
\newcommand{\F}{\mathcal F}
\title{Is the $5/2$ exponent optimal?\\Lower bounds and their exact scope}
\author{Fourier decomposition project: research proof draft}
\date{22 October 2026}
\begin{document}
\maketitle
\begin{abstract}
We prove an $\Omega(n^{5/2})$ gate lower bound for two explicitly restricted
architectures underlying the project's $\widetilde O(n^{5/2})$ symmetric-group
QFT constructions. The first uses separately implemented, fresh-input local
interlacing transforms; its local obstruction holds even at a small constant
error, for every sufficiently large diagram size. The second uses separately
implemented normalized-column blocks and bounded odd-polynomial amplification.
Neither result proves optimality for unrestricted QFT circuits. For those,
we give an elementary $\Omega(n\log n)$ bound and identify the missing
composition step. All coherent preprocessing and cleanup are charged.
These are research derivations with exact finite checks and internal audits,
not externally reviewed theorems or a claim of priority.
\end{abstract}

\section{Model and the distinction that matters}

Write $D=n!$. The Fourier transform convention is
\[
 \F_n\ket{g}=\sum_{\lambda\vdash n}\sum_{Q,R}
 \sqrt{d_\lambda/D}\,R_\lambda(g)_{QR}\ket{\lambda,Q,R},
\]
where $Q,R$ are standard tableaux of shape $\lambda$, $d_\lambda$ is their
number, and $R_\lambda$ is Young's orthogonal representation. We use injective
computational-basis encodings. An elementary gate acts on at most $a$ qubits,
for a fixed constant $a$. Ancillas start in zero. Clean error $\delta$ means
operator-norm error of the complete input isometry, including all output
workspace, relative to the ideal transform with zero workspace. The lower
bounds allow arbitrary bounded-arity unitaries and hence also hold for a
fixed finite gate set. There is no free quantum advice, data-dependent
preprocessing, QRAM operation, or unbounded-arity gate.

The word ``modular'' below is an explicit restriction on the circuit, not a
property inferred from the mathematical factorization of its target.
Correctness of the final product does not imply correctness of each factor.
Nor does it force a circuit to preserve the factors as separate gate sequences.

\section{A constant-error local lower bound}

For strictly interlacing integer addable and removable contents
$x_0>y_0>x_1>\cdots>y_{k-2}>x_{k-1}$, the first column of the local
orthogonal matrix has positive entries
\begin{equation}\label{eq:residue}
 \alpha_i^2=\frac{\prod_j|x_i-y_j|}
 {\prod_{h\ne i}|x_i-x_h|},\qquad \sum_i\alpha_i^2=1.
\end{equation}
Only this column is needed for the lower bound. Its input label is denoted
$\star$. For a partition $\lambda$, its canonical Frobenius record is
$d$ and the arms $A_i=\lambda_i-i$, legs $B_i=\lambda'_i-i$, $1\le i\le d$,
where $d$ is its Durfee-square side. Every word occupies a fixed physical
register. Output labels may be corner ordinals or added-row numbers.

\subsection{A same-size family with disjoint parameter blocks}

Let $m\ge1$, $K=4m$, and $s\in\{0,1\}^m$. Set
\begin{align*}
 x_i&=8m-4i &&(0\le i\le K),\\
 y_j&=x_j-c_{j\bmod4}+s_{\lfloor j/4\rfloor}e_{j\bmod4}
 &&(0\le j<K),\\
 c&=(2,2,1,3),& e&=(1,-1,-1,1).
\end{align*}
These strictly interlace and are the corners of genuine partitions.
For a direct construction, in block $r$ put $W_r=8(m-r)$ and use the eight
row lengths
\[
\begin{array}{ll}
s_r=0:&(W_r,W_r,W_r-2,W_r-2,W_r-4,W_r-7,W_r-7,W_r-7),\\
s_r=1:&(W_r,W_r-3,W_r-3,W_r-3,W_r-4,W_r-4,W_r-6,W_r-6).
\end{array}
\]
They concatenate to positive nonincreasing rows. Each block has area
$8W_r-29$, so every partition has
\[
 N_m=32m^2+3m
\]
boxes. Its corner lists are the displayed $x,y$. The selected corner
$x_{4r}$ is always the addable row $8r+1$. Its residue changes on a
zero-to-one flip of $s_r$ by
\begin{equation}\label{eq:flip}
 \frac{\alpha_{4r}(s_r=1)^2}{\alpha_{4r}(s_r=0)^2}
 =\frac12\frac76\frac{10}{9}\frac{14}{15}=\frac{49}{81}.
\end{equation}
This ratio is independent of all other bits.

Restrict variable bits to $1\le r<\lfloor m/2\rfloor$, fixing all others
at zero. All shapes have Durfee side $4m$. Changed cells lie strictly to
the right of column $4m$, so all legs are fixed; bit $s_r$ changes only
the eight arms in rows $8r+1,\ldots,8r+8$. These blocks are disjoint.
There are $b_m=\max(0,\lfloor m/2\rfloor-1)$ variable bits.

\begin{lemma}[Uniformly visible coordinates]\label{lem:mass}
For every bit string of the unpadded family and $0\le r<m$,
$\alpha_{4r}^2\ge1/(32m)$.
\end{lemma}
\begin{proof}
Compare with midpoint contents $y_j^{\rm mid}=x_j-2$. Their residue at
index $i$ is
\[
 b_i=\frac{\binom{2i}{i}}{4^i}
      \frac{\binom{2(K-i)}{K-i}}{4^{K-i}}\ge\frac1{2K}.
\]
Indeed $\binom{2j}{j}/4^j\ge1/(2\sqrt j)$ for $j\ge1$, by induction
using $(2j+1)^2\ge4j(j+1)$; the endpoints also satisfy the displayed bound.
At $i=4r$, block $b$ with $t=b-r$ multiplies the midpoint residue by
\[
 P_0(t)=1-\frac5{(16t+10)(16t+14)},\qquad
 P_1(t)=1-\frac5{(16t+2)(16t+6)}.
\]
Both are in $(0,1]$. At $t=0,-1$ each is at least $7/12$. On the
remaining rays $t=j$ and $t=-j-1$, $j\ge1$, each deficit is at most
$5/(256j^2)$. Since $\prod(1-u_j)\ge1-\sum u_j$ for $u_j\in[0,1]$,
and $\sum_{j\ge1}j^{-2}\le2$, the tail product is at least $59/64$.
Omitting any of these factors only raises the product. Consequently
\[
 \alpha_{4r}^2\ge\frac1{8m}\left(\frac7{12}\right)^2\frac{59}{64}
 >\frac1{32m}.
\]
\end{proof}

\subsection{Padding and the averaging argument}

For any $N\ge35$, choose the largest $m$ with $N_m\le N$ and add
$p=N-N_m$ boxes to the first row. Here $0\le p\le64m+34$, and
$m=\sqrt{N/32}+O(1)$. The frozen bit $s_0=0$ makes this padding common
to all records. It changes only the first Frobenius arm and preserves
the disjoint variable blocks. At a selected persistent corner $x_{4r}$,
$r\ge1$, put $d=8m-x_{4r}=16r$. Cancellation of the first-row corner
factors in \eqref{eq:residue} gives the multiplier
\begin{equation}\label{eq:pad}
 \frac{d(d+p-1)}{(d-1)(d+p)}
 =1+\frac{p}{(d-1)(d+p)}\ge1.
\end{equation}
Equivalently it is $\prod_{h=0}^{p-1}(d+h)^2/((d+h)^2-1)$.
It is independent of the bits, so both Lemma~\ref{lem:mass} and the
$7/9$ positive-amplitude flip survive. The output row remains $8r+1$;
its corner ordinal shifts by one if $p>0$, independently of the bits.

\begin{theorem}[Fresh-input constant-error obstruction]\label{thm:local}
Let $N\ge35$. Suppose a fixed circuit receives the canonical Frobenius record of an
$N$-box shape and the label $\star$, with all other inputs independent of
the shape. It must preserve the record and produce the positive column
\eqref{eq:residue}, with complete-output error at most $\delta$ for every
valid record. All its preprocessing and cleanup are charged. Let $m$ be
as above. If $u$ variable blocks in the padded family are not touched by
any gate, then
\[
 \delta^2\ge\frac{u}{2592m}.
\]
In particular, for $\delta\le1/144$, its gate count satisfies
\begin{equation}\label{eq:localcount}
 G_N\ge\left\lceil\frac{
 \max(0,\lfloor m/2\rfloor-1-\lfloor m/8\rfloor)}a\right\rceil
 =\Omega(\sqrt N).
\end{equation}
\end{theorem}
\begin{proof}
Let $U$ be the untouched blocks. Fix the bits outside $U$. The circuit
is the identity on the union of untouched physical registers. After
stripping their unchanged labels, its remaining output is independent of
all bits in $U$. Project onto the correct remaining parameter labels and
zero work; write $v_r$ for its amplitude at output row $8r+1$. This number
is independent of every bit in $U$, even if the actual circuit only
approximately preserves the accessed parameter registers.

For every choice of the untouched bits, orthogonality of the selected
output rows gives
\[
 \delta^2\ge\sum_{r\in U}|v_r-\alpha_{4r}(s)|^2.
\]
Average uniformly over these bits. Conditional on all bits other than
$s_r$, its two positive amplitudes are $h_r$ and $7h_r/9$, so their
variance is $h_r^2/81\ge1/(2592m)$. The mean squared error of any
constant is at least the variance. Summing proves the claim.
At the stated error $u\le\lfloor m/8\rfloor$. Every gate meets at most
$a$ disjoint blocks, which proves \eqref{eq:localcount}. For $m\ge16$
at least $m/4$ blocks are touched. The target column belongs to a matrix
of dimension $\Theta(m)$, so this is also an $\Omega(k)$ local bound.
\end{proof}

\section{Optimality within clean streaming modules}

\begin{definition}[Fresh full-promise modular streaming architecture]
For each outer size $M=2,\ldots,n$ and prefix size $N=0,\ldots,M-1$,
the circuit contains a separately counted local module. Each module has
the interface of Theorem~\ref{thm:local} (with trivial smaller-size cases),
is correct on every valid $N$-box record with error at most $1/144$, and
has its control fixed to one common active value independent of the record.
It cannot read an external tableau, redundant encoding,
or persistent shape-dependent cache. All preprocessing is inside the
module. Gates of different modules are not shared, fused, or cancelled.
\end{definition}

\begin{theorem}[Restricted $5/2$ optimality]\label{thm:stream}
Every circuit in this architecture uses $\Omega(n^{5/2})$ elementary
gates. More precisely, summing the integer lower bound
\eqref{eq:localcount} gives
\[
 G\ge\sum_{N=35}^{n-1}(n-N)G_N^{\rm lower}
 =\frac{n^{5/2}}{40\sqrt2\,a}+O(n^2).
\]
\end{theorem}
\begin{proof}
Each prefix size $N$ occurs at exactly $n-N$ outer sizes. Its forced
touched-block count is $3m/8+O(1)=3\sqrt N/(32\sqrt2)+O(1)$.
Rounding contributes only $O(n^2)$ to the sum. Finally
$\sum_{N<n}(n-N)\sqrt N=(4/15)n^{5/2}+O(n^{3/2})$.
\end{proof}

Different modules may have different witness shapes. This is legitimate
because each fixed module has a full-promise correctness contract; we do
not claim that one tableau path encounters all witnesses simultaneously.
The theorem matches the existing near-linear-core implementation's
$\widetilde O(n^{5/2})$ exponent, including coherent coefficient computation.
It does not apply once its explicit interface assumptions are removed.

\section{A second restricted obstruction: direct-column amplification}

Consider the alternative induction construction at size $M$. Its unitary
induction matrix is $F$. Normalized column preparation retains the input
tag $i\in\{1,\ldots,M\}$, and a fixed normalized tag-erasure bra has
coefficients $v_i$. Allow per-column attenuation $|c_i|\le1$.

\begin{lemma}[Normalization and polynomial degree]\label{lem:poly}
If this procedure gives the scalar block $F/\mu$, then $\mu\ge\sqrt M$.
If a subsequent scalar amplifier has successful block $p(1/\mu)F$,
where $p$ is an odd polynomial of degree at most $q$ and
$|p(x)|\le1$ for $x\in[-1,1]$, then target error $\varepsilon<1$ implies
\[
 q\ge\frac{1-\varepsilon}{\arcsin(1/\mu)}
 \ge\frac{2(1-\varepsilon)}\pi\sqrt M.
\]
\end{lemma}
\begin{proof}
Since $F$ is invertible, scalar normalization requires
$c_i\overline v_i=1/\mu$ for every $i$. Thus $|v_i|\ge1/\mu$ and
$1\ge\sum_i|v_i|^2\ge M/\mu^2$.
Put $T(t)=p(\sin t)$. It is a trigonometric polynomial of degree at most
$q$, bounded by one, with $T(0)=0$. Bernstein's inequality
\cite[Theorem 1.1]{Bernstein} gives $\norm{T'}_\infty\le q$.
Since $|T(\arcsin(1/\mu))|\ge1-\varepsilon$, integration proves the
first bound. The second uses $\arcsin x\le\pi x/2$ for $0\le x\le1$.
\end{proof}

This is a scalar odd-polynomial amplification model, as in the appropriate
QSVT construction \cite{QSVT}; it is not a query bound for every circuit
that knows the explicit matrix $F$.

\begin{lemma}[Physical cost of the prescribed normalized-column block]
Let $M\ge2$. Use the construction's fixed row-word tableau encoding, and
let $W_M$ be its coherent full-promise module for all valid incoming shapes
$\lambda\vdash M-1$, including $(M-1)$. Its prescribed action on a basis
source $(\lambda,i,P)$ has tableau factor $F_\lambda\ket{i,P}$ in every
attenuation branch, tensored with a normalized tag/flag state. This is
the full column action, not just a small successful block. A separate
implementation of $W_M$ with complete-output error $\eta<1/\sqrt2$
requires at least $\lceil(M-2)/a\rceil$ gates.
\end{lemma}
\begin{proof}
Take the incoming shape $\lambda=(M-1)$. Its induced representation is
the natural permutation representation on $e_1,\ldots,e_M$. A Young basis
is the uniform vector and
\[
 h_j=\frac{(j-1)e_j-\sum_{i<j}e_i}{\sqrt{j(j-1)}},\qquad 2\le j\le M.
\]
The branching state is $e_M$; the transversal $t_i=(i,\ldots,M)$ sends
it to $e_i$. Therefore on tag input $i=j$, the final tableau has second-row
label $j$ with probability $(j-1)/j\ge1/2$. The attenuation flag and tag
unpreparation act on other registers and preserve this tableau marginal.
For each $j=2,\ldots,M-1$, its input row slot initially stores row 1.
Choose a bit distinguishing row 1 from row 2 in that fixed slot. If it is
untouched, the actual state has zero projection onto the required event,
giving vector error at least $\sqrt{(j-1)/j}$. Thus all $M-2$ selected
physical bits must be touched. A gate touches at most $a$ of them.
\end{proof}

\begin{proposition}[Modular amplification architecture]
Suppose each outer induction independently satisfies the scalar-polynomial
contract of Lemma~\ref{lem:poly} at some fixed error $\varepsilon<1$.
Here $q$ counts calls in the ideal scalar-polynomial template; compiled
approximate calls need not produce an exactly polynomial block.
Suppose these $q$ calls to $W_M$ or $W_M^\dagger$ are nonoverlapping separately
implemented modules, each with the full incoming-shape promise: a forward
call implements the prescribed full column action, and an inverse call is
its inverse circuit or implements the full inverse unitary. Require
$\eta<1/\sqrt2$ for these module implementations.
Then this architecture uses $\Omega(n^{5/2})$ gates.
\end{proposition}
\begin{proof}
The bounds multiply under these stated module contracts, giving
$\Omega(\sqrt M)\Omega(M)=\Omega(M^{3/2})$ gates per outer size.
Sum over $M\le n$. Approximate full inverse implementations obey the same
physical bound by taking their adjoint.
\end{proof}

An arbitrary unitary completion having only successful block $F/\mu$
need not have the large tableau marginal used above. This is why the
physical lemma is stated for the prescribed full column action. Likewise,
the proposition does not cover fused calls, alternative encodings or
amplifiers, or directly synthesizing $F$ without this block encoding.

\section{What can be proved without those architecture restrictions?}

\begin{theorem}[Unrestricted support bound]\label{thm:general}
Every bounded-arity clean circuit for $\F_n$ at error $0\le\delta<1$
has
\[
 G\ge\frac{\log_2(n!)+2\log_2(1-\delta^2/2)}a
 =\Omega(n\log n).
\]
This holds for arbitrary connectivity, unlimited clean ancillas, and any
injective computational-basis encodings of the valid input/output labels.
\end{theorem}
\begin{proof}
The trivial-irrep label $\ket\tau$ is a computational basis state, and
$\F_n^\dagger\ket\tau=\ket u=D^{-1/2}\sum_g\ket g$.
The forward clean guarantee implies
$\norm{C^\dagger\ket{\tau,0}-\ket{u,0}}\le\delta$ by unitarity.
Starting from a computational basis state, each $a$-qubit gate multiplies
the number of supported computational basis strings by at most $2^a$.
Hence $C^\dagger\ket{\tau,0}$ has support at most $2^{aG}$.
A normalized state of support $s$ has overlap with $\ket{u,0}$ at most
$\sqrt{s/D}$, whereas vector error $\delta$ forces overlap at least
$1-\delta^2/2$. Rearranging proves the bound.
\end{proof}

There is also a bound that counts only gates crossing the right-tableau
register, when that register is physically separate. On identity input,
\[
 \F_n\ket e=\sum_{\lambda,Q}\sqrt{d_\lambda/D}
       \ket{\lambda,Q}\ket Q.
\]
Its Schmidt probabilities are $d_\lambda/D$, each repeated $d_\lambda$
times. Since $d_\lambda\le\sqrt D$, each is at most $D^{-1/2}$.
A rank-$R$ state has squared overlap at most $R/\sqrt D$ with this target.
Each crossing gate of total arity $a$ multiplies Schmidt rank by at most
$2^a$; put all ancillas on the other side. Thus
\[
 G_{\rm cross}\ge
 \frac{\tfrac12\log_2(n!)+2\log_2(1-\delta^2/2)}a.
\]
This remains valid with free gates acting entirely on either side of the
cut, but is still only
$\Omega(n\log n)$, not $\Omega(n^{5/2})$.

\section{Why the gap cannot be removed by multiplying local bounds}

A simple example exposes the missing inference. For a classical parameter
string $x$, define a controlled rotation
\[
 U_x=R_z\left(\theta\sum_{j=1}^k x_j\right).
\]
A standalone implementation at accuracy resolving the effect of one bit
must touch all $k$ parameter bits. Nevertheless
\[
 U_x^T=R_z\left(T\theta\sum_{j=1}^k x_j\right)
\]
can be synthesized by computing the Hamming weight once, applying one
rotation with the multiplied angle, and uncomputing. For $T$ and inverse
error polynomially bounded in $k$, this costs $\widetilde O(k)$, not
$\Omega(Tk)$.
One can choose $\theta=1/(Tk)$ and inverse-polynomial local error to
keep the individual bit changes resolvable without phase wrapping.
This is not a claimed shortcut for the symmetric-group QFT; it shows why
fresh-input bounds alone are not an unrestricted direct-sum theorem.

The missing result is a lower bound surviving shared shape-dependent
workspace, changing intermediate bases, and global resynthesis. Coherence
does not forbid any of these. Conversely, an improved algorithm could
exploit them, a different induction embedding, or a balanced subgroup merge.
Our results make the standalone-local-replacement strategy exponent-tight;
they do not establish that such other strategies are impossible.

Classical Fourier lower bounds must also be matched to the gate model.
Ailon's entropy result \cite{Ailon} charges a two-coordinate linear gate;
a physical qubit gate acts on many amplitude-coordinate pairs simultaneously.
It therefore does not supply the sought quantum gate bound. The
Cleve--Watrous lower bound \cite{CW} concerns depth for cyclic QFTs, not
an $n^{5/2}$ size lower bound for $S_n$.

\section{Verification and conclusion}

Exact-rational reference checks live in
\texttt{constant\_error\_local\_probe.py} and
\texttt{modular\_lower\_bound\_probe.py}; they verify family sizes,
corner residues, padding, disjoint Frobenius blocks, and integer accounting.
\texttt{amplification\_lower\_bound\_probe.py} checks the one-row induction
matrix and the amplification accounting. The uniform arguments above,
not finite numerical checks, are the proofs. The upper construction is
documented separately in \texttt{subcubic-qft-theorem.tex} and the research
log. The rigorous conclusion of this audit is:
\[
 \boxed{\begin{gathered}
 \text{$5/2$ is optimal for the stated modular architectures;}\\
 \text{unrestricted optimality is not proved.}
 \end{gathered}}
\]

\begin{thebibliography}{9}
\bibitem{Bernstein} H. Queff\'elec and R. Zarouf,
\emph{On Bernstein's inequality for polynomials}, Theorem 1.1.
\href{https://arxiv.org/abs/1903.10801}{arXiv:1903.10801}.
\bibitem{QSVT} A. Gily\'en, Y. Su, G. H. Low, and N. Wiebe,
\emph{Quantum singular value transformation and beyond}.
\href{https://arxiv.org/abs/1806.01838}{arXiv:1806.01838}.
\bibitem{Ailon} N. Ailon,
\emph{A Lower Bound for Fourier Transform Computation in a Linear Model
Over $2\times2$ Unitary Gates Using Matrix Entropy}.
\href{https://arxiv.org/abs/1305.4745}{arXiv:1305.4745}.
\bibitem{CW} R. Cleve and J. Watrous,
\emph{Fast parallel circuits for the quantum Fourier transform}.
\href{https://arxiv.org/abs/quant-ph/0006004}{arXiv:quant-ph/0006004}.
\end{thebibliography}
\end{document}
